\section{Proof of the Positive-Excess Lower Bound}\label{app:high-functional}

This appendix proves Lemma~\ref{lem:high-functional}. The proof has two
independent halves. The first half is a second-moment computation, which shows
that if a constant fraction of the starting edges carry a well spread family of
paths, then a long path survives the random deletion with constant
probability. The second half builds this family from the positive excess.

\medskip
\noindent\textbf{Notation of this appendix.} Besides the symbols of
Sections~\ref{subsec:notation} and~\ref{sec:difficult-graphs}, this appendix
uses the following notation: the set
\(\Pset{e}\) of candidate paths, the weights \(\wt{e}\), a path prefix \(\pref\),
the branching rate \(\brnum\), the fraction \(\gfrac\), and the drift
\(\drift\) and upper bound \(\ubd\) of the log-branching walk.

\subsection{Flows on the path tree}

Fix a directed edge \(e=(u,v)\), and let \(\Pset{e}\) be the set of paths of
length \(\plen\) that start at \(v\) and whose first step does not go to \(u\).
Give each \(P\in\Pset{e}\) a weight \(\wt{e}(P)\ge0\), and for a path prefix
\(\pref\) of length \(k\), write
\[
  \wt{e}(\pref):=\sum_{P\supseteq\pref}\wt{e}(P)
\]
for the total weight of the candidate paths that begin with \(\pref\).

Counting the candidate paths is not enough. If there are many of them and each
edge is kept with probability \(\keepp\), the expected number of kept paths may
be large, but the paths may share a long first segment, and then deleting one
common edge destroys all of them. We therefore need a family that is
\emph{spread out}, and we express this with a flow.

\begin{definition}[Flow]\label{def:flow}
Let \(\brnum>1\) and \(\gfrac\in(0,1]\). Let \(e=(u,v)\) be a directed edge such
that \(\ball{\plen}{v}\) is a tree. A weight function \(\wt{e}\) on \(\Pset{e}\)
is a \emph{\flowt{\brnum}{\gfrac}} if
\[
  \sum_{P\in\Pset{e}}\wt{e}(P)=1,
  \qquad
  \wt{e}(\pref)\le\frac{\brnum^{-k}}{\gfrac}
  \quad\text{for every prefix }\pref\text{ of length }k .
\]
If such a \(\wt{e}\) exists, we say that \emph{\(e\) carries a
\flowt{\brnum}{\gfrac}}.
\end{definition}

Thus a flow is a unit flow in which every prefix of length \(k\) has capacity
\(\brnum^{-k}/\gfrac\). On infinite trees, flows of this kind are used to define
the branching number (see Lyons~\cite{Lyons1990} and
Lyons--Peres~\cite{LyonsPeres2016}), but we do not need that theory, because the
flows we use are built directly in Lemma~\ref{lem:weight-construction}. The
next lemma is the standard second-moment argument which shows that percolation
survives when \(\keepp\brnum>1\).

\begin{lemma}[Deletion calculation for weighted paths]
\label{lem:weighted-path-percolation}
Suppose that \(\Gin\) has \(\nvert\) vertices, no self-loops, and
\(2\le\dg{\Gin}{v}\le\degb\). Let \(\brnum>1\) and \(\gfrac\in(0,1]\), and assume
that at least a \(\gfrac\) fraction of the directed edges carry a
\flowt{\brnum}{\gfrac}. If \(\keepp\brnum>1\), then
\begin{equation}\label{eq:abstract-path-lower-bound}
  \Lpath{\plen}{\Gin}{\keepp}
  \ge \frac{2\gfrac^2}{\degb}\left(1-(\keepp\brnum)^{-1}\right).
\end{equation}
\end{lemma}

This lemma does not use the excess, which is needed only in
Lemma~\ref{lem:weight-construction} to produce the flows.

\begin{proof}
First fix a directed edge \(e=(u,v)\) that carries a \flowt{\brnum}{\gfrac},
together with such a flow \(\wt{e}\). Keep each edge with probability
\(\keepp\), and define
\[
  Z_e:=\sum_{P\in\Pset{e}}\keepp^{-\plen}\wt{e}(P)
     \1_{\{\text{all edges of }P\text{ are kept}\}}.
\]
Each candidate path has \(\plen\) edges, so \(\E Z_e=\sum_P\wt{e}(P)=1\).

Take two candidate paths \(P,Q\), and let \(k(P,Q)\) be the length of their
common prefix. All candidate paths lie in the tree \(\ball{\plen}{v}\), where
two paths from \(v\) that separate never meet again, so \(P\) and \(Q\) have
exactly \(k(P,Q)\) common edges, and both are kept with probability
\(\keepp^{2\plen-k(P,Q)}\). Therefore
\begin{align*}
  \E Z_e^2
  &=\sum_{P,Q}\wt{e}(P)\wt{e}(Q)\keepp^{-k(P,Q)}\\
  &\le\sum_{k=0}^\plen \keepp^{-k}
       \sum_{|\pref|=k}\wt{e}(\pref)^2\\
  &\le\sum_{k=0}^\plen \keepp^{-k}
       \left(\max_{|\pref|=k}\wt{e}(\pref)\right)
       \sum_{|\pref|=k}\wt{e}(\pref)\\
  &\le \frac{1}{\gfrac}\sum_{k=0}^\plen(\keepp\brnum)^{-k}
   \le\frac{1}{\gfrac\bigl(1-(\keepp\brnum)^{-1}\bigr)}.
\end{align*}
For the second line, use
\(\keepp^{-k(P,Q)}\le\sum_{k=0}^{k(P,Q)}\keepp^{-k}\), and note that \(P\) and
\(Q\) have a common prefix of length \(k\) exactly when \(k\le k(P,Q)\); then
group the pairs \((P,Q)\) by this common prefix \(\pref\). For the last
line, note that every candidate path has exactly one prefix of each length
\(k\), so \(\sum_{|\pref|=k}\wt{e}(\pref)=1\). By the Cauchy--Schwarz
inequality,
\[
  \Prx{Z_e>0}\ge\frac{(\E Z_e)^2}{\E Z_e^2}
  \ge \gfrac\left(1-(\keepp\brnum)^{-1}\right),
\]
and if \(Z_e>0\), then some candidate path from \(v\) has all its edges kept.

Now we average over the directed edges. For \(e=(u,v)\), the event \(Z_e>0\) is
contained in the event that a kept path of length \(\plen\) starts at \(v\), and
a vertex \(v\) is the head of exactly \(\dg{\Gin}{v}\le\degb\) directed edges,
so
\[
  \sum_{e\text{ carrying a flow}}\Prx{Z_e>0}
  \le
  \degb\sum_{v\in\Vset}
  \Prx{\text{a kept path of length }\plen\text{ starts at }v}
  =\degb\,\nvert\,\Lpath{\plen}{\Gin}{\keepp}.
\]
Since the minimum degree is at least \(2\), the graph \(\Gin\) has at least
\(\nvert\) edges, that is, at least \(2\nvert\) directed edges, and at least
\(2\gfrac\nvert\) of them carry a flow. Hence
\[
  \Lpath{\plen}{\Gin}{\keepp}
  \ge \frac1{\degb\,\nvert}\cdot 2\gfrac\nvert\cdot
       \gfrac\left(1-(\keepp\brnum)^{-1}\right)
  =\frac{2\gfrac^2}{\degb}
       \left(1-(\keepp\brnum)^{-1}\right).
\]
\end{proof}

It remains to produce the flows, and this is where the positive excess is used.

\begin{lemma}[Construction of the flows]\label{lem:weight-construction}
There are constants \(\brnum=\brnum(\excd)\in\bigl(2,\,2(1+\excd)\bigr)\),
\(\gfrac=\gfrac(\degb,\excd)>0\) and \(\nontreez=\nontreez(\degb,\excd)>0\) with
the following property. Let \(\plen\ge1\) be an integer, and let \(\Gin\) satisfy
\Hyp{} and be \loctree{\plen}{\nontreez}. Then at least a \(\gfrac\) fraction of
the directed edges carry a \flowt{\brnum}{\gfrac}.
\end{lemma}

\begin{proof}[Proof of Lemma~\ref{lem:high-functional}]
Take \(\brnum,\gfrac,\nontreez\) from Lemma~\ref{lem:weight-construction}. Since
\(\brnum>2\), we have \(2/\brnum<1\), so we can fix
\[
  \keeppz:=\frac12\Bigl(\frac{2}{\brnum}+1\Bigr)\in\Bigl(\frac2\brnum,1\Bigr).
\]
Then \((\keeppz/2)\brnum>1\), and since Lemma~\ref{lem:weight-construction}
gives the hypothesis of Lemma~\ref{lem:weighted-path-percolation}, we can apply
that lemma with \(\keepp=\keeppz/2\):
\[
  \Lpath{\plen}{\Gin}{\keeppz/2}
  \ge\frac{2\gfrac^2}{\degb}
       \left(1-\left(\tfrac{\keeppz}{2}\brnum\right)^{-1}\right)=:\plow>0 .
\]
All of \(\brnum,\gfrac,\nontreez,\keeppz,\plow\) depend only on
\(\degb,\excd\), and none of them depends on \(\plen\), \(\nvert\) or \(\Gin\).
This proves Lemma~\ref{lem:high-functional}.
\end{proof}

Here we need \(2/\brnum<1\), that is, \(\brnum>2\), and we can choose such a
\(\brnum\) below \(2(1+\excd)\) exactly because \(\excd>0\);
see~\eqref{eq:lambda-choice} below.

\subsection{Constructing the flows}

\begin{proof}[Proof of Lemma~\ref{lem:weight-construction}]
\emph{The walk.} We use an auxiliary random walk only to define the weights;
the graph \(\Gin\) is fixed, and there is no edge deletion here. We write
\(\Prob_{\mathrm{nb}}\) for the law of the following nonbacktracking walk.

Choose a directed edge \(\vec e_0=(v_{-1},v_0)\) uniformly from all directed
edges of \(\Gin\). After reaching \(\vec e_i=(v_{i-1},v_i)\), choose
\(\vec e_{i+1}\) uniformly among the \(\dg{\Gin}{v_i}-1\) directed edges that
leave \(v_i\) and are not the reversal of \(\vec e_i\), and write
\(\vec e_{i+1}=(v_i,v_{i+1})\). Since every vertex has degree at least \(2\),
the walk never gets stuck.

Note that the walk forbids the reversal of the edge just used, not every edge
back to \(v_{i-1}\). The two rules differ only when \(v_{i-1}\) and \(v_i\) are
joined by parallel edges, and with the second rule such a step would have
fewer than \(\dg{\Gin}{v_i}-1\) choices, so the stationarity below would fail.
We also need that \(\Gin\) has no self-loops, since a self-loop at \(v_i\)
would give a step with \(v_{i+1}=v_i\).

For a fixed directed edge \(e\), we write \(\Prx[\mathrm{nb}]{\,\cdot\mid e}\)
for the law of the walk started with \(\vec e_0=e\).

\emph{Stationarity.} Fix a directed edge \((a,b)\). The directed edges that can
move to \((a,b)\) are the \(\dg{\Gin}{a}-1\) edges that enter \(a\), except
\((b,a)\), and each of them moves to \((a,b)\) with probability
\(1/(\dg{\Gin}{a}-1)\). Hence if \(\vec e_i\) is uniform over all directed
edges, then
\[
  \Prx[\mathrm{nb}]{\vec e_{i+1}=(a,b)}
  =(\dg{\Gin}{a}-1)\cdot\frac{1}{2|E(\Gin)|}\cdot\frac{1}{\dg{\Gin}{a}-1}
  =\frac{1}{2|E(\Gin)|},
\]
and \(\vec e_{i+1}\) is uniform as well. Since \(\vec e_0\) is uniform, it
follows that every \(\vec e_i\) is uniform, and since the same rule is used at
every step, for all \(t\ge0\) and \(j\ge1\) the blocks
\[
  (\vec e_0,\ldots,\vec e_{j-1})
  \quad\text{and}\quad
  (\vec e_t,\ldots,\vec e_{t+j-1})
\]
have the same joint law. We call this property \emph{stationarity}; note that
it does not mean that the steps are independent.

\emph{The weights.} Let \(\brnum>2\) be the constant fixed
in~\eqref{eq:lambda-choice} below. Call a walk of length \(\plen\) \emph{good}
if
\begin{equation}\label{eq:path-prefix-growth}
  \prod_{i=0}^{k-1}(\dg{\Gin}{v_i}-1)\ge\brnum^k
  \qquad\text{for every }1\le k\le \plen.
\end{equation}
In words, every prefix of a good walk has enough branching. Let \(\mathcal G\)
be the event that the walk is good, let
\(g(e):=\Prx[\mathrm{nb}]{\mathcal G\mid e}\), and if \(g(e)>0\), define
\[
  \wt{e}(P):=
  \frac{\Prx[\mathrm{nb}]{P\mid e}}{g(e)}
  \1_{\{P\in\mathcal G\}}.
\]
Thus \(\wt{e}\) is the law of the walk, restricted to good walks and
normalized.

Now let \(e=(v_{-1},v_0)\) be a directed edge such that \(\ball{\plen}{v_0}\) is
a tree. Then the walks of length \(\plen\) started with \(e\) are exactly the
paths in \(\Pset{e}\). Indeed, in a tree there is only one edge between
\(v_{-1}\) and \(v_0\), so forbidding this edge is the same as forbidding the
vertex \(v_{-1}\), and in a tree a nonbacktracking walk never comes back to a
vertex. Hence \(\wt{e}\) is a weight function on \(\Pset{e}\) with total mass
\(1\). To check the capacity bound, let \(\pref\) be a prefix of length \(k\).
If \(\pref\) is not the beginning of a good walk, then \(\wt{e}(\pref)=0\);
otherwise~\eqref{eq:path-prefix-growth} holds for \(\pref\), and
\begin{equation}\label{eq:finite-flow-bound}
  \wt{e}(\pref)
  \le \frac{\Prx[\mathrm{nb}]{\pref\mid e}}{g(e)}
  =\frac{1}{g(e)}\prod_{i=0}^{k-1}\frac{1}{\dg{\Gin}{v_i}-1}
  \le \frac{\brnum^{-k}}{g(e)} .
\end{equation}
Therefore, if \(g(e)\ge\gfrac\), then \(e\) carries a \flowt{\brnum}{\gfrac},
and it remains to find \(\brnum\) and \(\gfrac\) such that at least a \(\gfrac\)
fraction of the directed edges \(e=(u,v)\) have \(g(e)\ge\gfrac\) and a tree
ball \(\ball{\plen}{v}\).

\emph{Choice of \(\brnum\).} Let \(d:=2|E(\Gin)|/\nvert\) be the average degree,
so that \(d\ge3+2\excd\) by \Hyp{}. Since \(\vec e_i\) is uniform, the vertex
\(v_i\) is chosen with probability proportional to its degree:
\[
  \Prx[\mathrm{nb}]{v_i=v}=\frac{\dg{\Gin}{v}}{2|E(\Gin)|}.
\]
The function \(f(s)=s\log(s-1)\) is convex for \(s\ge2\), since
\(f''(s)=(s-2)/(s-1)^2\ge0\). By Jensen's inequality,
\begin{equation}\label{eq:entropy-gap}
  \E_{\mathrm{nb}}\log(\dg{\Gin}{v_i}-1)
  =\frac{\sum_v f(\dg{\Gin}{v})}{\nvert d}
  \ge\frac{f(d)}{d}
  =\log(d-1)
  \ge\log\bigl(2(1+\excd)\bigr).
\end{equation}
Fix
\begin{equation}\label{eq:lambda-choice}
  \brnum=\brnum(\excd):=2(1+\excd)^{1/4}\in\bigl(2,\ 2(1+\excd)\bigr).
\end{equation}
This interval is nonempty because \(\excd>0\), and the bound
\(d\ge3+2\excd\) in~\eqref{eq:entropy-gap} and this choice are the only places
where the positive excess is used. The upper end \(2(1+\excd)\) makes the drift
below positive, while the lower end \(2\) leaves room for the retention
probability \(\keeppz/2\).

\emph{The drift.} Define
\[
  Y_i:=\log(\dg{\Gin}{v_i}-1)-\log\brnum ,
\]
so that \(\sum_{i<k}Y_i\ge0\) is the same as the inequality
in~\eqref{eq:path-prefix-growth}. Since \(Y_i\) is a function of
\(\vec e_i\), the sequence \((Y_i)_{i\ge0}\) is stationary, and
by~\eqref{eq:entropy-gap},
\[
  \E_{\mathrm{nb}} Y_i\ \ge\ \log\bigl(2(1+\excd)\bigr)-\log\brnum
  \ =:\ \drift\ >\ 0,
  \qquad
  Y_i\le\log(\degb-1)=:\ubd .
\]

We need more than a positive average, because every prefix sum up to \(\plen\)
must be nonnegative. The following lemma, proved in
Section~\ref{subsec:positive-prefix}, gives this.

\begin{lemma}[Finite positive-prefix lemma]\label{lem:positive-prefix}
Let \((Y_i)_{i\ge0}\) be a stationary sequence of real random variables, that
is, for all integers \(t\ge0\) and \(j\ge1\) the blocks
\((Y_0,\ldots,Y_{j-1})\) and \((Y_t,\ldots,Y_{t+j-1})\) have the same joint
law. Assume that
\[
  Y_i\le \ubd,
  \qquad
  \E Y_i\ge\drift>0 .
\]
Then there is a constant \(\gfrac\in(0,1/4]\), depending only on \(\ubd\) and
\(\drift\), such that for every integer \(\plen\ge1\),
\begin{equation}\label{eq:positive-prefix}
  \Prx{\sum_{i=0}^{k-1}Y_i\ge0\text{ for every }1\le k\le \plen}
  \ge 4\gfrac.
\end{equation}
\end{lemma}

\emph{Many good starting edges.} By Lemma~\ref{lem:positive-prefix}, the walk
is good with probability at least \(4\gfrac\), and since the starting edge is
uniform,
\[
  \frac{1}{2|E(\Gin)|}\sum_{e}g(e)\ \ge\ 4\gfrac ,
\]
where the sum is over all directed edges. Since \(0\le g(e)\le1\), at least a
\(2\gfrac\) fraction of the directed edges satisfy
\begin{equation}\label{eq:good-directed-edge}
  g(e)\ge2\gfrac,
\end{equation}
since otherwise the average would be less than \(2\gfrac+2\gfrac=4\gfrac\).

\emph{The tree condition.} Choose \(\nontreez>0\) with
\(\degb\nontreez/2\le\gfrac\). Since \(\Gin\) is \loctree{\plen}{\nontreez}, at
most \(\nontreez\nvert\) vertices have a non-tree ball of radius \(\plen\), and
at most \(\degb\nontreez\nvert\) directed edges end at these vertices. There are
at least \(2\nvert\) directed edges in total, so at most a \(\gfrac\) fraction
of them end at such a vertex. Hence at least a \(\gfrac\) fraction of the
directed edges \(e=(u,v)\) satisfy~\eqref{eq:good-directed-edge} and have a
tree ball \(\ball{\plen}{v}\), and by~\eqref{eq:finite-flow-bound} each of them
carries a \flowt{\brnum}{\gfrac}. Since the constants \(\brnum\), \(\gfrac\)
and \(\nontreez\) depend only on \(\degb,\excd\), this proves
Lemma~\ref{lem:weight-construction}.
\end{proof}


\subsection{Proof of the finite positive-prefix lemma}
\label{subsec:positive-prefix}

\begin{proof}[Proof of Lemma~\ref{lem:positive-prefix}]
The two assumptions give \(\drift\le\E Y_0\le \ubd\).

\emph{A deterministic fact.} Let \(y_0,\ldots,y_{N-1}\le \ubd\) be real numbers
with \(W:=\sum_{i=0}^{N-1}y_i>0\), and extend the sequence periodically by
\(y_{i+N}=y_i\). Set \(s_k:=\sum_{i=0}^{k-1}y_i\), so that \(s_0=0\) and
\(s_{k+N}=s_k+W\), and call a position \(i\in\{0,\ldots,N-1\}\) \emph{good} if
\[
  s_{i+k}-s_i=\sum_{j=i}^{i+k-1}y_j\ge0
  \qquad\text{for every }1\le k\le N.
\]
We show that at least \(W/\ubd\) positions are good. Let
\[
  \underline{s}_i:=\min_{j\ge i}s_j
\]
be the smallest partial sum from position \(i\) on; the minimum exists because
the partial sums grow by \(W>0\) in every period. Since
\(\underline{s}_i=\min\{s_i,\underline{s}_{i+1}\}\), we have
\(\underline{s}_{i+1}\ge \underline{s}_i\). If
\(\underline{s}_{i+1}>\underline{s}_i\), then \(\underline{s}_i=s_i\), hence
\(\underline{s}_{i+1}\le s_{i+1}=s_i+y_i\le\underline{s}_i+\ubd\). Also
\(s_{i+k}\ge \underline{s}_{i+1}>s_i\) for every \(k\ge1\), so position \(i\) is
good. Thus every step \(\underline{s}_{i+1}-\underline{s}_i\) is either \(0\)
or at most \(\ubd\), and it is positive only at a good position. Finally,
\[
  \sum_{i=0}^{N-1}(\underline{s}_{i+1}-\underline{s}_i)
  =\underline{s}_N-\underline{s}_0=W,
\]
because \(\underline{s}_N=\min_{j\ge0}(s_j+W)=\underline{s}_0+W\). Hence at
least \(W/\ubd\) of these steps are positive, and at least \(W/\ubd\) positions
are good.

\emph{A block with positive sum.} Fix \(\plen\), let
\(N:=\lceil 4\ubd/\drift\rceil\,\plen\), and set
\(W_N:=\sum_{i=0}^{N-1}Y_i\). By stationarity,
\[
  \E W_N\ge\drift N,
  \qquad
  W_N\le\ubd N.
\]
Therefore,
\begin{equation}\label{eq:block-positive}
  \Prx{W_N\ge\frac\drift2N}
  \ge\frac{\drift}{2\ubd}.
\end{equation}
Indeed, if the probability on the left is \(a\), then
\(\drift N\le\E W_N\le a\ubd N+(1-a)\drift N/2\), which gives
\(a\ge\drift/(2\ubd-\drift)\ge\drift/(2\ubd)\).

\emph{Many good positions.} On the event in~\eqref{eq:block-positive}, the
deterministic fact with \(y_i=Y_i\) gives at least \(\drift N/(2\ubd)\) good
positions. A good position \(i\) with \(i+\plen\le N\) gives
\(\sum_{j=i}^{i+k-1}Y_j\ge0\) for all \(1\le k\le\plen\), using only terms of
the block and no periodic copies, and at most \(\plen\) positions have
\(i+\plen>N\). Since \(N\ge 4\ubd\plen/\drift\), we have
\(\plen\le\drift N/(4\ubd)\), so at least
\[
  \frac{\drift N}{2\ubd}-\plen\ge\frac{\drift N}{4\ubd}
\]
good positions have \(i+\plen\le N\).

\emph{Back to position \(0\).} Let \(J\) be uniform on \(\{0,\ldots,N-1\}\) and
independent of the sequence, and let \(\mathcal E_i\) be the event that
\(\sum_{j=i}^{i+k-1}Y_j\ge0\) for every \(1\le k\le \plen\). On the event
in~\eqref{eq:block-positive}, at least a \(\drift/(4\ubd)\) fraction of the
positions \(i\) satisfy \(\mathcal E_i\), so
\[
  \Prx{\mathcal E_J}
  \ge\frac{\drift}{4\ubd}\,
    \Prx{W_N\ge\frac\drift2N}
  \ge\frac{\drift^2}{8\ubd^2}.
\]
By stationarity, \(\Prx{\mathcal E_i}=\Prx{\mathcal E_0}\) for every \(i\), and
therefore \(\Prx{\mathcal E_J}=\Prx{\mathcal E_0}\). If we set
\[
  \gfrac:=\frac{\drift^2}{32\ubd^2},
\]
then \(\Prx{\mathcal E_0}\ge4\gfrac\), and \(\drift\le \ubd\) gives
\(\gfrac\le1/32\le1/4\). The constant \(\gfrac\) depends only on \(\ubd\) and
\(\drift\), and not on \(\plen\).
\end{proof}
